\section{Introduction}
\label{sec:intro}

Neuro-symbolic legal reasoning pipelines couple an LLM, which extracts structured predicates from a fact pattern, with a symbolic rule engine that applies an auditable rule to reach a decision~\cite{yang2025neuro}. Such systems are attractive because the rule structure is fixed and inspectable, in principle allowing the decision path to be verified independently of the model that produced it. We instantiate this division of labour as PYTHEN, a neuro-symbolic pipeline for hearsay classification under LegalBench~\cite{guha2023legalbench}, and show that it improves F1-score and reduces instability across model families relative to both direct and chain-of-thought prompting. An oracle experiment, in which gold predicates are fed directly into the rule tree, further shows that the symbolic component reaches perfect agreement with ground truth once its inputs are correct - indicating that essentially all residual error of the pipeline originates upstream, in the LLM's predicate extraction step.

This result exposes a gap that standard evaluation practice does not close. Predicate extraction is typically assessed by aggregate precision/recall/F1~\cite{olek2023evaluationf1scorerecent}, which implicitly treats every predicate error as equally consequential. This assumption fails once predicates are consumed by a rule tree: a rule tree encodes conditional structure, so some predicates sit on a decisive branch where a single wrong value flips the final verdict, while others are redundant with correlated conditions, short-circuited by an earlier branch, or on a path the specific instance never reaches. Two models with identical predicate-level F1 can therefore differ substantially in how many of their errors actually reach the final decision, and aggregate extraction metrics cannot distinguish the two cases, because they are computed independently of the rule structure the predicates feed into.

This paper asks, for each predicate error an LLM extractor makes, not only \emph{that} it occurred but \emph{whether it matters} - whether correcting it would change the system's output. We introduce a predicate-level intervention procedure: for each fact pattern, we correct one erroneous predicate at a time toward its gold value, rerun the PYTHEN rule tree, and observe whether the decision changes. This yields a Decision Impact Rate (DIR) for every predicate, which we compare against standalone extraction F1, and use to refine an extraction-error taxonomy along a decision-relevance dimension. Across three LLMs, we find that extraction difficulty and decision impact diverge sharply - the hardest predicate to extract is also the most decision-critical, while a comparably hard predicate is almost always inconsequential - and that nominally identical error types carry markedly different downstream risk across models. These results argue for rule-aware error attribution, rather than aggregate predicate accuracy alone, as a necessary complement when evaluating and verifying neuro-symbolic legal reasoning pipelines.
